\documentclass[12pt,a4paper]{article}

\usepackage{amsmath,amssymb}
\usepackage{hyperref}
\usepackage{natbib}
\usepackage{geometry}
\usepackage{booktabs}
\usepackage{xcolor}
\usepackage{microtype}

\geometry{margin=1in}

\hypersetup{
  colorlinks=true,
  linkcolor=blue,
  citecolor=blue,
  urlcolor=blue
}

% ─────────────────────────────────────────────────────────────────────────────
\title{\textbf{Non-Locality and the Boundary of Reality:\\
Quantum Potential, Irreversible Actualization,\\
and the Thermodynamic Origin of Classical Existence}}

\author{Heath W.\ Mahaffey\\
\small Independent Researcher, Entiat, Washington 98822, USA\\
\small \href{mailto:hmahaffeyges@gmail.com}{hmahaffeyges@gmail.com}\\
\small \url{https://github.com/hmahaffeyges/IAM-Validation}}

\date{March 2026\\
\small Zenodo DOI: \href{https://doi.org/10.5281/zenodo.18702042}{10.5281/zenodo.18702042}}
% ─────────────────────────────────────────────────────────────────────────────

\begin{document}
\maketitle

% ─────────────────────────────────────────────────────────────────────────────
\begin{abstract}
Quantum non-locality --- the apparent ability of entangled systems to
exhibit correlated behavior regardless of spatial separation --- has been
called the defining mystery of quantum mechanics. We argue that it is
not mysterious at all, but is the expected behavior of physical systems
that have not yet crossed the boundary between potential and actual.
Within the Informational Actualization Model (IAM), that boundary has
a precise thermodynamic identity: every irreversible quantum-to-classical
transition pays a cost of $k_B T_H \ln 2$ per bit to the cosmic horizon,
where $T_H = \hbar H / 2\pi k_B$ is the Gibbons--Hawking temperature.
The crossing is physically real, thermodynamically irreversible, and
cosmologically consequential.

Before this crossing, a quantum system is not a particle with hidden
properties. It is a field of potential --- a superposition that has not
yet been written to the ledger of classical existence. Non-locality is
the expected property of potential: potential has no position, no
trajectory, no classical identity. After the crossing, the system
becomes a classical particle, the information is encoded permanently
on the cosmic horizon, and locality is restored. There is no spooky
action at a distance because there is no particle to act at a distance
until actualization occurs.

This identification --- that the boundary of classical existence is a
thermodynamic threshold, not a mathematical postulate --- is a consequence
of the Jacobson--Cai-Kim derivation extended to include the Landauer cost
of irreversible bulk information production. The accumulated ledger of
13.8 billion years of such crossings is measured by the activation
function $\mathcal{E}(a) = \exp(1 - 1/a)$, derived from horizon
thermodynamics. Its cosmological consequences --- $\mu(a) < 1$ for matter
on timelike worldlines, $\Sigma(a) = 1$ for photons on null geodesics ---
are confirmed against Planck 2018 through 17 converged MCMC chains with
no free parameters beyond $\Lambda$CDM. The bridge from the quantum
boundary to the cosmological scale is the critical exponent $n = 5/2$,
derived independently from quantum decoherence physics (bottom-up) and
from horizon thermodynamic consistency (top-down). Their convergence
constitutes a derivation, not a scaling argument.

The predicted matter-sector modification today is $\mu_0 - 1 = -0.136$,
derived from first principles with zero free parameters.
\end{abstract}

% ─────────────────────────────────────────────────────────────────────────────
\section{The Mystery and the Category Error}
\label{sec:intro}
% ─────────────────────────────────────────────────────────────────────────────

Quantum non-locality has been presented as one of the deepest mysteries
in physics. Entangled particles, separated by arbitrary distances,
exhibit correlated measurement outcomes that cannot be explained by any
local hidden variable theory~\citep{Bell1964}. The correlations violate
Bell inequalities~\citep{Aspect1982}, the violation has been confirmed to
high confidence~\citep{Hensen2015}, and the conclusion seems inescapable:
nature is non-local.

We propose that this conclusion rests on a category error. The mystery of
non-locality arises from applying the language of classical particles ---
position, trajectory, definite state --- to physical systems that are not
yet classical particles. Before a quantum-to-classical transition occurs,
there is no particle. There is a quantum field in superposition: a
distribution of potential that has not paid the thermodynamic cost of
becoming actual. Asking where the entangled ``particle'' is, or how it
communicates with its partner, is like asking where a wave is before it
breaks. The question assumes a classical object that does not yet exist.

Within the Informational Actualization Model
(IAM)~\citep{IAM_theory,IAM_obs}, this distinction is not philosophical.
It is derived. The boundary between potential and actual has a precise
thermodynamic identity: every irreversible quantum-to-classical transition
pays a Landauer cost~\citep{Landauer1961} of $k_B T_H \ln 2$ per bit to
the cosmic horizon. Before payment, the system is potential. After
payment, it is classical. The boundary is real, irreversible, and
cosmologically consequential.

% ─────────────────────────────────────────────────────────────────────────────
\section{The Ontological Distinction: Potential and Actual}
\label{sec:ontology}
% ─────────────────────────────────────────────────────────────────────────────

Standard quantum mechanics describes two regimes that are treated as
limits of a single formalism: the quantum regime of superposition and
entanglement, and the classical regime of definite states and particles.
The transition between them is handled by the measurement postulate ---
a postulate, not a derivation. IAM derives it.

The distinction is ontological, not merely mathematical. A quantum system
in superposition is not a classical particle with unknown properties. It
is categorically different: a field of potential governed by the
Schr\"{o}dinger equation, with no definite position, no definite
momentum, and no classical identity. The word ``particle'' already
assumes actualization. Applied to a pre-measurement quantum system, it
is a category error that has generated a century of apparent paradoxes.

IAM identifies the crossing from potential to actual with a specific
physical event: an irreversible thermodynamic interaction between a
quantum system and matter on a timelike worldline that dissipates
energy $Q \geq k_B T \ln 2$ --- the Landauer threshold~\citep{Landauer1961}.
This is not a postulate. It follows from requiring that the entropy
budget of the Jacobson--Cai-Kim horizon thermodynamic
derivation~\citep{Jacobson1995,CaiKim2005} be complete: every bit of
classical information produced by an irreversible quantum-to-classical
transition must be encoded somewhere. The only available surface at
cosmological scales is the cosmic horizon.

The sector separation follows immediately. Photons travel null geodesics:
$d\tau = 0$. They do not accumulate proper time, do not participate in
gravitational decoherence, and do not write to the cosmic ledger. They
remain in the potential sector indefinitely. Massive particles travel
timelike worldlines, participate in gravitational decoherence, and pay
the Landauer cost at each virialization event. They are in the matter
sector~\citep{IAM_time}.

This sector separation is not imposed. It is derived from the geometry
of the manifold combined with the thermodynamic requirement that
classical information production carry a cost. It should be noted that
local environmental decoherence --- from lab walls, air molecules, or
phonon coupling --- is not a separate mechanism that bypasses the cosmic
ledger. Every irreversible quantum-to-classical transition, local or
cosmic, writes to the horizon. Local environments are relays, not
substitutes: they mediate the irreversible interaction that produces the
information, but the information itself is encoded on the cosmic horizon
as the ultimate boundary of classical existence. The sector split governs
which systems participate in this process at all, not where the
information is first absorbed.

% ─────────────────────────────────────────────────────────────────────────────
\section{The Derivation Chain: From Jacobson to the Cosmic Ledger}
\label{sec:derivation}
% ─────────────────────────────────────────────────────────────────────────────

The derivation chain is explicit and traceable. It proceeds in five steps.

\textbf{Step 1: Jacobson (1995).} The Einstein field equations are an
equation of state derived from the proportionality of entropy and area
for local causal horizons~\citep{Jacobson1995}. Gravity is thermodynamic.
This derivation assumes the entropy functional contains only geometric
terms.

\textbf{Step 2: Cai and Kim (2005).} The Friedmann equations follow from
the same thermodynamic argument applied to the FRW apparent
horizon~\citep{CaiKim2005}. Both derivations assume the entropy functional
is complete.

\textbf{Step 3: The missing contribution.} Gravitational structure
formation is the dominant source of irreversible classical information
production in the universe. Each virialization event converts quantum
superpositions to classical outcomes. The virial theorem for $1/r$
potentials partitions the gravitational binding energy exactly:
\begin{equation}
2K = |V|, \qquad K = \tfrac{1}{2}|V|.
\label{eq:virial}
\end{equation}
The kinetic half $K$ is the decoherence energy --- the Landauer cost of
each quantum-to-classical transition. This identification carries a
thermodynamic cost of $k_B T_H \ln 2$ per bit, where
$T_H = \hbar H / 2\pi k_B$ is the Gibbons--Hawking
temperature~\citep{Gibbons1977}. The information must be encoded on the
cosmic horizon.

\textbf{Step 4: The accumulated ledger.} Integrating the information
production rate across 13.8 billion years of gravitational structure
formation yields the activation function:
\begin{equation}
\mathcal{E}(a) = \exp\!\left(1 - \frac{1}{a}\right),
\label{eq:Ea}
\end{equation}
derived from the requirement that the cumulative informational entropy be
consistent with the Jacobson--Cai-Kim horizon thermodynamics. This
function is monotonically increasing by construction ($d\mathcal{E}/da >
0$ always), begins at the electroweak transition ($\mathcal{E} \to 0$ as
$a \to 0$), equals unity today ($\mathcal{E}(1) = 1$), and asymptotes to
$e$ without arriving ($\mathcal{E}(\infty) = e$)~\citep{IAM_time}.

\textbf{Step 5: The cosmological consequence.} Including the informational
entropy contribution in the Cai-Kim first law modifies the matter
perturbation equation through an additional Hubble friction term. The
result is the dual-sector modification:
\begin{equation}
\mu(a) = \frac{H^2_{\Lambda\mathrm{CDM}}}{H^2_{\Lambda\mathrm{CDM}}
+ \beta_m\,\mathcal{E}(a)} < 1, \qquad \Sigma(a) = 1,
\label{eq:sector}
\end{equation}
where $\beta_m = \Omega_m/2$ is derived from the virial theorem with zero
free parameters~\citep{IAM_theory}. The Planck 2018 MCMC posterior
returns $\beta_m = 0.1583 \pm 0.0033$, consistent with the virial
prediction at $0.2\sigma$.

The derivation is complete. Every step traces from an established result.
Nothing is postulated beyond what Jacobson and Cai-Kim already assumed,
plus the requirement that the entropy budget be complete.

% ─────────────────────────────────────────────────────────────────────────────
\section{What Non-Locality Actually Is}
\label{sec:nonlocality}
% ─────────────────────────────────────────────────────────────────────────────

With the ontological distinction established and the derivation chain
explicit, the nature of quantum non-locality becomes clear.

Entangled systems are fields of potential that have not yet paid the
Landauer cost of actualization. They have no classical identity, no
definite position, no trajectory. They are not particles separated by a
distance. They are one undivided field of potential that has not yet been
partitioned into classical outcomes by an irreversible thermodynamic
event.

Non-locality is the expected property of potential. Potential is not
local. Locality is a property of actual things --- classical particles with
definite positions on timelike worldlines. Before actualization, the
concept of spatial separation does not apply in the classical sense,
because there is no classical object to be separated.

When measurement occurs --- when an irreversible Landauer transaction is
forced on a timelike worldline --- the field of potential actualizes. At
that moment it becomes a particle, writes to the cosmic ledger, and
locality is restored. There is no mechanism required for the apparent
instantaneous correlation because there is no signal traveling between
two separated particles. There is one field of potential becoming
classical at one end, and the entanglement resolving as a consequence.

Bell's theorem~\citep{Bell1964} is precisely correct: there are no local
hidden variables. IAM agrees, but for a derived reason rather than a
postulated one. There are no hidden variables because there is no
classical particle to have hidden variables until actualization occurs.
Bell's result is exactly what one would predict from IAM's ontology.

The sector assignment makes a quantitative prediction. Photon-based
entangled systems ($\Sigma = 1$) do not decohere during transit. Their
entanglement is preserved at arbitrary distances, and the CHSH
$S$-parameter remains $|S| = 2\sqrt{2}$ identically, in full agreement
with quantum mechanics. Matter-based entangled systems ($\mu < 1$)
experience gravitational decoherence on timelike worldlines. The Bell
violation decays as~\citep{IAM_QM}:
\begin{equation}
S(t) = 2\sqrt{2}\,(1 - D(t)), \qquad
D(t) = \mathcal{E}(t/\tau_{\mathrm{IAM}})/e,
\label{eq:bell}
\end{equation}
where $\tau_{\mathrm{IAM}}$ is the gravitational decoherence timescale,
derived from first principles~\citep{IAM_grav}:
\begin{equation}
\tau_{\mathrm{IAM}} = \frac{\hbar k_B^2 T^2 \ln 2}{E_G^3},
\label{eq:tauIAM}
\end{equation}
where $E_G = Gm^2/r$ is the gravitational self-energy of the
superposition and $T$ is the bath temperature. This formula contains
no free parameters. Table~\ref{tab:decoherence} gives representative
values across three experimental regimes and compares IAM against
Penrose--Di\'{o}si and standard environmental decoherence.

\begin{table}[h]
\centering
\small
\caption{Gravitational decoherence timescales across three experimental
regimes. IAM predicts a ramp profile with $T^2$ temperature scaling and
$m^{-6}$ mass scaling. Penrose--Di\'{o}si predicts exponential decay
with no temperature dependence and $m^{-5/3}$ scaling. The mesoscopic
frontier ($\sim 10^{-12}$~kg) is where IAM and Penrose--Di\'{o}si
predictions diverge measurably with current experimental reach.}
\label{tab:decoherence}
\begin{tabular}{p{4.2cm}p{2.8cm}p{2.8cm}p{3.6cm}}
\toprule
\textbf{System} & \textbf{IAM $\tau_{\mathrm{IAM}}$} &
\textbf{P--D $\tau_{\rm PD}$} & \textbf{Key discriminator} \\
\midrule
Trapped ion ($\sim 10^{-25}$~kg) &
  $\gg$ practical & $\gg$ practical & Both irrelevant \\
Superconducting qubit ($\sim 10^{-20}$~kg) &
  $\gg$ practical & $\gg$ practical & Both irrelevant \\
Nanosphere ($10^{-12}$~kg, 10~mK) &
  $\approx 560~\mu$s & $\sim 10~\mu$s &
  Profile shape; $T^2$ vs $T^0$ \\
\bottomrule
\end{tabular}
\end{table}

The photon-matter asymmetry --- photon entanglement maintained at
1,200~km (Micius satellite, 2017~\citep{Yin2017}) versus matter
entanglement limited to approximately 1.3~m (trapped ions,
2019~\citep{Hensen2015}) --- is a ratio of one million, consistent with
the sector separation.\footnote{Recent results have demonstrated Bell
inequality violations via quantum indistinguishability through path
identity, without requiring entanglement in the conventional
sense~\citep{SciAdv2025}. This does not challenge the IAM framework:
such violations arise from indistinguishability structure, not from
entangled fields of potential in the IAM sense. The sector-dependent
decoherence prediction applies specifically to systems whose correlations
originate from a shared quantum state, not from post-selection on
indistinguishable paths.}

A clarification on massive unstable particles is warranted. Entanglement
has been confirmed between top quark pairs at the Large Hadron
Collider~\citep{CMS2024}, despite top quarks being massive fermions on
timelike worldlines. This is consistent with IAM. Gravitational
decoherence operates on virialization timescales --- the timescale for
gravitational collapse to convert quantum superpositions to classical
outcomes, which is orders of magnitude longer than the top quark lifetime
of $\sim 10^{-25}$~s. Top quarks decay before gravitational decoherence
can act. The IAM prediction is that entanglement between massive systems
decays on \textit{gravitational} decoherence timescales, not on particle
decay timescales. The two mechanisms are distinct.

A further open question is entanglement across time rather than space.
Experiments have demonstrated correlations between particles that never
coexisted simultaneously~\citep{Megidish2013}, suggesting entanglement
is a property of the quantum state rather than of spatially separated
objects. In IAM, the field of potential exists prior to actualization
regardless of whether the actualization events are spacelike or timelike
separated. The temporal case is not excluded by the framework, but a
full treatment of temporal entanglement within IAM is beyond the scope
of this paper and is identified as an open direction.

% ─────────────────────────────────────────────────────────────────────────────
\section{The Bridge: $n = 5/2$ from Quantum to Cosmic}
\label{sec:bridge}
% ─────────────────────────────────────────────────────────────────────────────

The claim that a single thermodynamic mechanism operates from the
quantum boundary to the cosmological scale requires a bridge. That
bridge is the critical exponent $n = 5/2$, derived independently from
two directions.

\subsection{Bottom-Up: From Quantum Decoherence Physics}

The Joos--Zeh result establishes that gravitational decoherence of
macroscopic collapsed structures is instantaneous on cosmological
timescales~\citep{JoosZeh1985}. The rate-limiting step is therefore
virialization, not decoherence itself. Integrating the information
production rate over the Press--Schechter halo mass
function~\citep{PressSchechter1974}, with peak height scaling
$\nu \propto D(a)^{-1/2}$ for halos near the characteristic mass $M_*$:
\begin{equation}
\dot{I} \propto \rho_m(a)\cdot D(a)^{5/2} \cdot f(a) \cdot H(a),
\label{eq:Idot_bottomup}
\end{equation}
with $n = 5/2$ derived from quantum Darwinism applied to gravitational
structure formation~\citep{IAM_Zurek}.

\subsection{Top-Down: From Horizon Thermodynamic Consistency}

The top-down derivation requires the cumulative informational entropy to
reproduce $\mathcal{E}(a) = \exp(1-1/a)$. Substituting
matter-domination scalings and writing $\dot{I} \propto \rho_m D^n H$
(the derivation is presented under matter-domination as the epoch
during which the bulk of structure formation occurs; the full treatment
across all epochs is given in~\citep{IAM_Zurek}):
\begin{equation}
S_{\rm info}(a) \propto \int a^{n - 9/2 - 1}\,da \propto a^{n-9/2}.
\end{equation}
Requiring $S_{\rm info}(a) \propto a^{-1}$ (from the $\exp(1-1/a)$
functional form) forces:
\begin{equation}
n - \frac{9}{2} = -1 \quad\Longrightarrow\quad \boxed{n = \frac{5}{2}}.
\label{eq:n_topdown}
\end{equation}

\subsection{Convergence}

Both directions yield $n = 5/2$. The bottom-up derivation obtains this
from quantum decoherence physics integrated over the Press--Schechter
mass function. The top-down derivation obtains it from requiring the
activation function to be consistent with horizon thermodynamics. Their
convergence is not a coincidence. It constitutes a derivation connecting
the quantum boundary of actualization to cosmological observables from
two independent directions simultaneously~\citep{IAM_Zurek}.

The same irreversible thermodynamic process that resolves entanglement at
the quantum scale accumulates into the ledger $\mathcal{E}(a)$ that
suppresses structure growth at the cosmological scale. One mechanism. One
law.

% ─────────────────────────────────────────────────────────────────────────────
\section{IAM's Law}
\label{sec:law}
% ─────────────────────────────────────────────────────────────────────────────

The results of the preceding sections can be stated as a single
thermodynamic law:

\medskip
\noindent\fbox{\parbox{0.93\textwidth}{
\textbf{IAM's Law.} Every irreversible crossing from potential to actual
--- every quantum-to-classical transition on a timelike worldline ---
pays a thermodynamic cost of $k_B T_H \ln 2$ per bit to the cosmic
horizon. This cost is physically real, thermodynamically irreversible,
and cannot be recovered. Classical existence is continuously purchased,
one decoherence event at a time.
}}
\medskip

This law has three direct consequences:

\textbf{1. Particles are post-actualization concepts.} Before the
Landauer cost is paid, there is no particle. There is a field of
potential. The word ``particle'' applied to a pre-measurement quantum
system is a category error.

\textbf{2. Non-locality is the expected behavior of potential.}
Entangled systems are fields of potential that have not paid the cost of
actualization. Non-locality is not a mystery to be explained --- it is
the natural state of potential before the boundary is crossed.

\textbf{3. The arrow of time is a thermodynamic necessity.} The cosmic
ledger $\mathcal{E}(a)$ is monotonically increasing because each
Landauer transaction is irreversible. Time has direction because
actualization has direction. The ledger only grows~\citep{IAM_time}.

% ─────────────────────────────────────────────────────────────────────────────
\section{Observational Confirmation}
\label{sec:obs}
% ─────────────────────────────────────────────────────────────────────────────

IAM's law is not a philosophical claim. It is a physical claim with
measured consequences.

Seventeen MCMC chains were run against the full Planck 2018 likelihood
(TT, TE, EE + low-ell + lensing). All 17 converged to $R-1 < 0.01$
with no discarded runs, no parameter adjustments, and no excluded
redshift ranges. The results~\citep{IAM_obs}:

\begin{itemize}
\item $\sigma_8 = 0.7998 \pm 0.0058$ (suppressed from $\Lambda$CDM's
  $0.809$, consistent with KiDS-Legacy, DES Y3, HSC Y3 at $0.1\sigma$)
\item $H_0^{\rm (photon)} = 67.16 \pm 0.47$ km\,s$^{-1}$\,Mpc$^{-1}$
  ($0.37\sigma$ from Planck)
\item $H_0^{\rm (matter)} = 72.26$ km\,s$^{-1}$\,Mpc$^{-1}$
  ($0.75\sigma$ from SH0ES)
\item $\beta_m = 0.1583 \pm 0.0033$ ($0.2\sigma$ from virial prediction
  $0.1575$)
\item $\Delta\chi^2 = +0.54$ best-fit vs $\Lambda$CDM
\end{itemize}

The sector-dependent Hubble split --- photon-sector observables on null
geodesics vs matter-sector observables on timelike worldlines --- is a
direct consequence of IAM's law. The two ``values'' of $H_0$ are not
a tension. They are measurements of two distinct physical quantities
separated by the accumulated Landauer cost of 13.8 billion years of
actualization.

The predicted matter-sector modification today is $\mu_0 - 1 = -0.136$,
derived from the virial theorem coupling $\beta_m = \Omega_m/2$ and the
activation function $\mathcal{E}(a)$, with zero free parameters beyond
$\Lambda$CDM. The existing chain results are already consistent with this
value. The Euclid DR1 measurement (October 2026), with projected
sensitivity $\sigma(\mu_0) \approx 0.04$, will determine the precision
of the mechanism's magnitude --- not whether the mechanism operates, which
the Planck 2018 chains have already established.

% ─────────────────────────────────────────────────────────────────────────────
\section{Discussion}
\label{sec:discussion}
% ─────────────────────────────────────────────────────────────────────────────

The framework presented here does not modify quantum mechanics. It
identifies what quantum mechanics has always been describing: the
behavior of potential before it crosses the thermodynamic boundary into
classical existence.

Standard quantum mechanics is the correct theory of potential. General
relativity is the correct theory of actual spacetime geometry. IAM is
the thermodynamic account of the boundary between them --- the cost of
crossing, the ledger of crossings, and the cosmological consequence of
13.8 billion years of such crossings accumulated on the cosmic horizon.

The measurement problem, the arrow of time, quantum non-locality, and
the Hubble and $\sigma_8$ tensions are not separate problems. They are
different aspects of one problem: the incomplete entropy budget of the
Jacobson--Cai-Kim derivation. Including the Landauer cost of
irreversible bulk information production closes the budget, addresses
the foundational paradoxes, and produces testable cosmological
predictions.

The fact that decoherence is thermodynamically irreversible has been
accepted since the 1980s~\citep{JoosZeh1985,Zurek1981}. The fact that
horizons carry entropy has been accepted since
Bekenstein~\citep{Bekenstein1973} and Hawking~\citep{Hawking1975}. The
fact that gravity is thermodynamic has been accepted since
Jacobson~\citep{Jacobson1995}. The missing step was asking where the
Landauer cost of every irreversible quantum-to-classical transition
goes at cosmological scales, and what happens when 13.8 billion years
of such costs accumulate.

The answer is $\mathcal{E}(a)$. The consequences are $\mu < 1$ and
$\Sigma = 1$. The confirmation is 17 converged chains against the full
Planck 2018 likelihood.

% ─────────────────────────────────────────────────────────────────────────────
\section*{Acknowledgments}
% ─────────────────────────────────────────────────────────────────────────────

The thermodynamic framework of T.\ Jacobson~\citep{Jacobson1995} is the
foundation upon which this work is built. His 1995 derivation of the
Einstein equations as an equation of state established that gravity is
thermodynamic; IAM is the completion of that derivation when the entropy
budget is made explicit. The quantum Darwinism framework of
W.H.\ Zurek~\citep{Zurek2003,Zurek2009} provided the local account of
the quantum-to-classical transition that this paper extends to
cosmological scales. The thermal time program of C.\
Rovelli~\citep{Connes1994} posed the foundational question of the
thermodynamic origin of time that the activation function
$\mathcal{E}(a)$ addresses. A personal response from
Professor P.J.E.\ Peebles identified the quantum-to-classical derivation
as the foundational gap requiring closure before the cosmological
framework could be trusted --- that identification directly motivated the
bridge from quantum decoherence physics to the cosmic ledger presented
in Section~\ref{sec:bridge}.

% ─────────────────────────────────────────────────────────────────────────────
\section*{Data Availability}
% ─────────────────────────────────────────────────────────────────────────────

All MCMC chains, analysis scripts, and prediction figures are publicly
archived at \url{https://github.com/hmahaffeyges/IAM-Validation} with
canonical Zenodo DOI
\href{https://doi.org/10.5281/zenodo.18702042}{10.5281/zenodo.18702042}.
Repository timestamps confirm all predictions were established prior to
any observational comparison.

\bibliographystyle{plainnat}
\bibliography{iam_nonlocality_boundary}

\end{document}
